\section{Estratégias de crescimento}
\label{sec:crescimento}

\subsection{Critérios de escolha}

A orientação sugere considerar expansão de mercado, lançamento de novos produtos ou serviços, parcerias estratégicas e estratégias de marketing e aquisição. A matriz de produto e mercado de \citeonline{ansoff1957} distingue quatro direções de crescimento: penetração de mercado, desenvolvimento de mercado, desenvolvimento de produto e diversificação. Dado que a \projeto{} ainda não penetrou nenhum mercado e opera com recursos escassos, a diversificação foi descartada. As demais opções foram avaliadas segundo quatro critérios de viabilidade, apresentados na Tabela~\ref{tab:opcoes}.

\begin{table}[htb]
\centering
\caption{Avaliação das opções de crescimento}
\label{tab:opcoes}
\fontsizetable
\begin{tabularx}{\textwidth}{L{2.7cm} L{2.0cm} L{2.2cm} L{2.2cm} Y}
\toprule
\textbf{Opção} & \textbf{Custo inicial} & \textbf{Aderência à fase de validação} & \textbf{Exigência da equipe} & \textbf{Efeito na economia unitária}\\
\midrule
Parcerias com clínicas e ONGs & Baixo & Alta & Média & Positivo: cada clínica paga R\$~79,90 por mês e traz atendimentos.\\
Expansão geográfica (nova praça) & Médio & Média & Alta & Neutro até haver densidade; risco de diluir esforço.\\
Marketing e aquisição de tutores & Baixo a médio & Média & Média & Depende do CAC, ainda não medido.\\
Novos produtos e serviços & Médio & Baixa a média & Alta & Positivo só se aumentar assinatura ou reduzir o fundo.\\
\bottomrule
\end{tabularx}
\fonteautores
\end{table}

Com base nessa avaliação, definem-se duas estratégias que se reforçam mutuamente: a Estratégia 1, de densidade local por meio de parcerias (principal), e a Estratégia 2, de aprofundamento de receita e novos serviços (complementar).

\subsection{Estratégia 1 (principal): densidade local por meio de parcerias}

\textbf{Descrição.} Em plataformas de dois lados, o valor para cada lado depende da presença do outro, e o arranque exige conquistar ambos \cite{rochet2003}. \citeonline{eisenmann2006} recomendam decidir qual lado subsidiar e por quanto tempo. Na \projeto{}, o lado subsidiado é o tutor, cujo núcleo é gratuito, e o lado pagante inclui as clínicas (Plano Parceiro) e os patrocinadores. A estratégia consiste em concentrar o esforço em Belo Horizonte e em sua região metropolitana, construindo densidade de clínicas e tutores em poucas áreas antes de abrir novas praças. A capital integra o grupo das dez cidades com mais startups do país, o que favorece o acesso a parceiros e a talentos \cite{diario2025bh}.

\textbf{Componentes.}
\begin{alineas}
  \item \textit{Parcerias}: três clínicas no piloto, com três meses gratuitos do Plano Parceiro, e organizações não governamentais e lideranças comunitárias como porta de entrada dos tutores, com a faculdade de medicina veterinária como possível apoio à curadoria da triagem;
  \item \textit{Aquisição}: os três canais prioritários de comunicação e de distribuição da Atividade VIII \cite{grupo2026a8}, a saber, redes sociais, parcerias com clínicas e indicação, e código QR no balcão das clínicas e lojas oficiais de aplicativos, com programa de indicação remunerado em créditos;
  \item \textit{Expansão por praça}: abrir uma nova área apenas quando os gatilhos da Tabela~\ref{tab:gatilhos} forem atingidos.
\end{alineas}

\begin{table}[htb]
\centering
\caption{Gatilhos para abrir uma nova praça}
\label{tab:gatilhos}
\fontsizetable
\begin{tabularx}{\textwidth}{L{5.2cm} Y}
\toprule
\textbf{Gatilho} & \textbf{Referência}\\
\midrule
Pelo menos 10 clínicas ativas e 1.500 usuários ativos mensais na praça atual & Premissas do mês 12 da Atividade IX\\
Custo de aquisição observado igual ou inferior ao teto & Subseção~\ref{sub:cac}\\
Pelo menos 5\% dos usuários assinando o PetCuida+ & Premissa da Atividade IX\\
Fundo de liquidação por atendimento igual ou inferior a R\$~6,60, com trajetória de queda & Tabela~\ref{tab:alavancas}\\
Convênio ou patrocínio contratado & Dependência do resultado positivo\\
\bottomrule
\end{tabularx}
\fonteautores
\end{table}

\subsection{Estratégia 2 (complementar): aprofundamento de receita e novos serviços}

\textbf{Descrição.} Em termos da matriz de Ansoff, trata-se de desenvolvimento de produto para o mesmo mercado. A estratégia atua sobre as três fontes de receita já definidas e sobre o principal ponto fraco do modelo, o fundo de liquidação. Seus componentes são: (i) realizar o teste de aceitação do preço e elevar a conversão ao PetCuida+ de 5\% para 8\% dos usuários; (ii) vender campanhas patrocinadas que financiem parte do fundo, objetivo de reduzi-lo de R\$~6,60 para R\$~3,30 por atendimento; (iii) ampliar o Plano Parceiro, com agenda e destaque no mapa; e (iv) acrescentar novos serviços somente quando houver evidência de demanda.

\textbf{Novos serviços candidatos.} Os testes de usabilidade registraram sugestões espontâneas de adoção, controle de alimentação e comparação de preços em lojas parceiras \cite{grupo2026a6}. Esses itens entram em lista de espera e só são desenvolvidos se, após a validação do núcleo, houver demanda recorrente de usuários e se o custo de desenvolvimento for coberto por receita ou patrocínio. Essa regra evita a dispersão de uma equipe pequena, risco apontado na Atividade XI.

\subsection{Relação entre as estratégias}

A Estratégia 1 gera os usuários e as clínicas sem os quais a Estratégia 2 não tem base de conversão; a Estratégia 2 gera a margem que torna sustentável a expansão da Estratégia 1. A sequência é, portanto, simultânea e interdependente: o piloto (Estratégia 1) fornece os dados para o teste de preço e para a venda de patrocínio (Estratégia 2).
